\section{模型假设}
在题目给定参数与参考排程口径下，采用以下简化假设：
\begin{enumerate}
  \item 货箱作为不可拆分任务处理，其质量、体积、类别和时限按输入记录取值；本文不对输入数据作统计插补。
  \item 同一运输架次按给定闭环路线完成服务，架次开始、返航及逐箱送达时刻采用参考结果记录的时间口径。
  \item 无人机在一次运输任务中连续占用，电池从任务开始占用至达到下一次可用时刻；充电周转由结果文件中的可用时刻表示。
  \item 通信模式和中继任务沿用参考结果，不在本文中重新搜索候选位置或重优化通信链路。
  \item 问题四冻结运输和通信任务时序；各分区独立执行且不跨组共享无人机、电池或中继资源。
  \item 峰值资源需求按任务区间并发数计量。该指标表示固定时序下的并发配置下界，不等同于实体设备 ID 使用数量，也不单独刻画组间转场。
\end{enumerate}
上述假设用于明确参考结果的适用范围；风雨扰动、临时禁飞区、链路随机衰落和资源故障不在本稿数值范围内。
